\begin{abstract}
Segmentation models in medical imaging are trained on one acquisition setting
and must generalize to others. Label-generative domain randomization, of which
SynthSeg is the archetype, addresses this by discarding the acquired scan and
resampling intensities from a per-label Gaussian mixture over a \emph{dense}
anatomical label map -- an annotation far richer than the segmentation task
itself provides, and one that does not preserve texture: real gradients and
partial-volume detail become piecewise noise. We ask when that loss costs
accuracy, and find the answer is a property of the \emph{task}. Where the
target is separated from its surroundings by an anatomical interface --
encapsulated abdominal organs, bounded by fat planes -- destroying texture
costs nothing. Where no interface exists -- diffuse glioma, infiltrating beyond
both the enhancing margin and the peritumoral FLAIR abnormality -- the target
must be inferred from tissue appearance, and destroying texture is expensive.
We exploit this with PALETTE, a single-source augmentation needing no dense
label map: it partitions a scan's foreground by unsupervised $k$-means,
sub-divides each class spatially by a Voronoi partition, and applies a signed
random affine remap to the \emph{real} intensities of each sub-region. Remapping
existing intensities rather than replacing them leaves edges, gradients and
partial-volume detail largely intact; only the contrast relationships between
regions are randomized. A controlled ablation isolates this: holding the
partition and per-region statistics fixed and swapping only the real-intensity
remap for label-style noise fill costs $7.7$ and $7.2$ Dice on the two tasks
with no tissue interface, against $+1.1$ and $-1.1$ on the two with one.
Across cross-contrast and cross-modality transfer on healthy-brain, MS-lesion,
brain-tumour and abdominal-organ segmentation, one nnU-Net configuration
trained with PALETTE improves significantly over all five competing methods, on
both Dice and boundary distance.
\end{abstract}
